\documentclass[10pt,a4paper]{amsart}
\usepackage{amsmath,amssymb,amsthm,graphicx}
% Retained minimal preamble for the Stacks Project chapter more-algebra.tex (More on Algebra).
% The chapter file pinned-original-source.tex carries only \input{preamble} on its first line, so no
% preamble text lives inside the chapter itself.  The Stacks Project uses one single shared preamble for
% every chapter; the snapshot stages it identically as sources/stacks-chow/preamble.tex and
% sources/stacks-etale-cohomology/preamble.tex.  The block below is copied mechanically, line for line,
% from sources/stacks-chow/preamble.tex lines 14-18 (the Xy-pic setup) and 154-217 (the native theorem
% declarations together with the mathematical macros).  Nothing was added, removed or redefined; the
% wrapper class and the amsmath/amssymb/amsthm/graphicx base packages are supplied by the assembler.
% Chapter: more-algebra.tex -- More on Algebra.
% For commutative diagrams we use Xy-pic
\usepackage[all]{xy}

% We use 2cell for 2-commutative diagrams.
\xyoption{2cell}

% Theorem environments.
%
\theoremstyle{plain}
\newtheorem{theorem}[subsection]{Theorem}
\newtheorem{proposition}[subsection]{Proposition}
\newtheorem{lemma}[subsection]{Lemma}

\theoremstyle{definition}
\newtheorem{definition}[subsection]{Definition}
\newtheorem{example}[subsection]{Example}
\newtheorem{exercise}[subsection]{Exercise}
\newtheorem{situation}[subsection]{Situation}

\theoremstyle{remark}
\newtheorem{remark}[subsection]{Remark}
\newtheorem{remarks}[subsection]{Remarks}

\numberwithin{equation}{subsection}

% Macros
%
\def\lim{\mathop{\mathrm{lim}}\nolimits}
\def\colim{\mathop{\mathrm{colim}}\nolimits}
\def\Spec{\mathop{\mathrm{Spec}}}
\def\Hom{\mathop{\mathrm{Hom}}\nolimits}
\def\Ext{\mathop{\mathrm{Ext}}\nolimits}
\def\SheafHom{\mathop{\mathcal{H}\!\mathit{om}}\nolimits}
\def\SheafExt{\mathop{\mathcal{E}\!\mathit{xt}}\nolimits}
\def\Sch{\mathit{Sch}}
\def\Mor{\mathop{\mathrm{Mor}}\nolimits}
\def\Ob{\mathop{\mathrm{Ob}}\nolimits}
\def\Sh{\mathop{\mathit{Sh}}\nolimits}
\def\NL{\mathop{N\!L}\nolimits}
\def\CH{\mathop{\mathrm{CH}}\nolimits}
\def\proetale{{pro\text{-}\acute{e}tale}}
\def\etale{{\acute{e}tale}}
\def\QCoh{\mathit{QCoh}}
\def\Ker{\mathop{\mathrm{Ker}}}
\def\Im{\mathop{\mathrm{Im}}}
\def\Coker{\mathop{\mathrm{Coker}}}
\def\Coim{\mathop{\mathrm{Coim}}}

% Boxtimes
%
\DeclareMathSymbol{\boxtimes}{\mathbin}{AMSa}{"02}

%
% Macros for moduli stacks/spaces
%
\def\QCohstack{\mathcal{QC}\!\mathit{oh}}
\def\Cohstack{\mathcal{C}\!\mathit{oh}}
\def\Spacesstack{\mathcal{S}\!\mathit{paces}}
\def\Quotfunctor{\mathrm{Quot}}
\def\Hilbfunctor{\mathrm{Hilb}}
\def\Curvesstack{\mathcal{C}\!\mathit{urves}}
\def\Polarizedstack{\mathcal{P}\!\mathit{olarized}}
\def\Complexesstack{\mathcal{C}\!\mathit{omplexes}}
% \Pic is the operator that assigns to X its picard group, usage \Pic(X)
% \Picardstack_{X/B} denotes the Picard stack of X over B
% \Picardfunctor_{X/B} denotes the Picard functor of X over B
\def\Pic{\mathop{\mathrm{Pic}}\nolimits}
\def\Picardstack{\mathcal{P}\!\mathit{ic}}
\def\Picardfunctor{\mathrm{Pic}}
\def\Deformationcategory{\mathcal{D}\!\mathit{ef}}
\usepackage{cleveref}
\crefname{originallemma}{original lemma}{original lemmas}
\Crefname{originallemma}{Original lemma}{Original lemmas}
\makeatletter
\newlabel{more-algebra-lemma-flatten-on-affine-blowup}{{053I}{source}{}{sourcepointer-1}{}}
\newlabel{more-algebra-lemma-flatten-on-affine-blowup@cref}{{[originallemma][0][]053I}{[1][0][]source}{}{}{}}
\makeatother
\crefname{originallemma}{original lemma}{original lemmas}
\Crefname{originallemma}{Original lemma}{Original lemmas}
\makeatletter
\newlabel{algebra-lemma-finitely-presented-over-subring}{{0561}{source}{}{sourcepointer-2}{}}
\newlabel{algebra-lemma-finitely-presented-over-subring@cref}{{[originallemma][0][]0561}{[1][0][]source}{}{}{}}
\makeatother
\crefname{originallemma}{original lemma}{original lemmas}
\Crefname{originallemma}{Original lemma}{Original lemmas}
\makeatletter
\newlabel{lemma-torsion}{{0537}{source}{}{sourcepointer-3}{}}
\newlabel{lemma-torsion@cref}{{[originallemma][0][]0537}{[1][0][]source}{}{}{}}
\makeatother
\crefname{originallemma}{original lemma}{original lemmas}
\Crefname{originallemma}{Original lemma}{Original lemmas}
\makeatletter
\newlabel{lemma-valuation-ring-torsion-free-flat}{{0539}{source}{}{sourcepointer-4}{}}
\newlabel{lemma-valuation-ring-torsion-free-flat@cref}{{[originallemma][0][]0539}{[1][0][]source}{}{}{}}
\makeatother
\crefname{originallemma}{original lemma}{original lemmas}
\Crefname{originallemma}{Original lemma}{Original lemmas}
\makeatletter
\newlabel{lemma-flat-finite-type-valuation-ring-finite-presentation}{{053E}{source}{}{sourcepointer-5}{}}
\newlabel{lemma-flat-finite-type-valuation-ring-finite-presentation@cref}{{[originallemma][0][]053E}{[1][0][]source}{}{}{}}
\makeatother
\crefname{originallemma}{original lemma}{original lemmas}
\Crefname{originallemma}{Original lemma}{Original lemmas}
\makeatletter
\newlabel{algebra-lemma-valuation-ring-colimit-affine-blowups}{{052M}{source}{}{sourcepointer-6}{}}
\newlabel{algebra-lemma-valuation-ring-colimit-affine-blowups@cref}{{[originallemma][0][]052M}{[1][0][]source}{}{}{}}
\makeatother
\crefname{originallemma}{original lemma}{original lemmas}
\Crefname{originallemma}{Original lemma}{Original lemmas}
\makeatletter
\newlabel{algebra-lemma-flat-finite-presentation-limit-flat}{{02JO}{source}{}{sourcepointer-7}{}}
\newlabel{algebra-lemma-flat-finite-presentation-limit-flat@cref}{{[originallemma][0][]02JO}{[1][0][]source}{}{}{}}
\makeatother
\crefname{originallemma}{original lemma}{original lemmas}
\Crefname{originallemma}{Original lemma}{Original lemmas}
\makeatletter
\newlabel{lemma-flat-torsion-free}{{0538}{source}{}{sourcepointer-8}{}}
\newlabel{lemma-flat-torsion-free@cref}{{[originallemma][0][]0538}{[1][0][]source}{}{}{}}
\makeatother
\title{Flattening a finite module along an affine blowup}
\author{The Stacks Project Authors}
\date{stacks-project:more-algebra.tex:local-snapshot-20261012; GNU FDL 1.2 or later}
\begin{document}
\maketitle
\setcounter{subsection}{1}
\noindent\textit{Editorial note.} All delivered mathematical text is the author's own, copied from the pinned Stacks Project chapter file more-algebra.tex; no statement or proof step was rewritten, reordered or removed.\par\medskip
\noindent\textit{Editorial note.} The chapter file pinned-original-source.tex carries only \textbackslash{}input\{preamble\} on its first line, so no preamble text lives inside the chapter. The retained file sources/stacks-more-algebra/retained-preamble.tex is a mechanical, line-for-line copy of the project's single shared preamble as staged in the snapshot (Xy-pic setup plus theorem declarations and mathematical macros). Nothing was added, removed or redefined, and the excerpt serves only as preamble, never as a span coordinate.\par\medskip
\noindent\textit{Editorial note.} Both counted spans start at the \textbackslash{}begin line of their environment, so the statement and the author's proof are delivered whole; surrounding prose of the section is deliberately outside the counted spans and is carried as standard prerequisites instead.\par\medskip
\noindent\textit{Editorial note.} The citation fields (source\_id, source\_version, authors, source\_url, license.name, source\_title) are identical to those of every other card this lane has produced from this chapter, so the batch selector treats them as one source.\par\medskip
\noindent\textit{Allowed standard prerequisites (editorial).}\begin{itemize}
\item Flatness of a finite module over a ring detected through the Fitting ideal of a presentation. the criterion counted by the sibling card S1-algebra-fitting-ideal at 1355-1411. Supplies the form in which flatness is verified at the end of the proof.
\item Affine blowup of a ring along an ideal, given by adjoining the ratios of the generators. used at proof lines 6097-6130. Provides the new ring over which the module is shown to be flat.
\item The equations of the blowup and their relation to the original presentation. used at proof lines 6130-6171. Supplies the columns added to the presentation matrix and the computation that follows.
\end{itemize}\medskip
\ifdefined\hypertarget\hypertarget{sourcepointer-1}{}\fi
\noindent\textit{Source reference.} Original lemma 053I: focused-seeds/AGP/sources/stacks-snapshot-20261011/tags.txt, source lines 6606--6606. The lemma's own canonical tag, quoted so that readers can resolve the label inside the delivered text; it points at the counted statement itself and carries no proof content.\par\medskip
\ifdefined\hypertarget\hypertarget{sourcepointer-2}{}\fi
\noindent\textit{Source reference.} Original lemma 0561: focused-seeds/AGP/sources/stacks-snapshot-20261011/tags.txt, source lines 6694--6694. Navigation pointer to a statement the counted proof refers to but does not contain; nothing load-bearing is delegated to it.\par\medskip
\ifdefined\hypertarget\hypertarget{sourcepointer-3}{}\fi
\noindent\textit{Source reference.} Original lemma 0537: focused-seeds/AGP/sources/stacks-snapshot-20261011/tags.txt, source lines 6595--6595. Navigation pointer to a statement the counted proof refers to but does not contain; nothing load-bearing is delegated to it.\par\medskip
\ifdefined\hypertarget\hypertarget{sourcepointer-4}{}\fi
\noindent\textit{Source reference.} Original lemma 0539: focused-seeds/AGP/sources/stacks-snapshot-20261011/tags.txt, source lines 6597--6597. Navigation pointer to a statement the counted proof refers to but does not contain; nothing load-bearing is delegated to it.\par\medskip
\ifdefined\hypertarget\hypertarget{sourcepointer-5}{}\fi
\noindent\textit{Source reference.} Original lemma 053E: focused-seeds/AGP/sources/stacks-snapshot-20261011/tags.txt, source lines 6602--6602. Navigation pointer to a statement the counted proof refers to but does not contain; nothing load-bearing is delegated to it.\par\medskip
\ifdefined\hypertarget\hypertarget{sourcepointer-6}{}\fi
\noindent\textit{Source reference.} Original lemma 052M: focused-seeds/AGP/sources/stacks-snapshot-20261011/tags.txt, source lines 6575--6575. Navigation pointer to a statement the counted proof refers to but does not contain; nothing load-bearing is delegated to it.\par\medskip
\ifdefined\hypertarget\hypertarget{sourcepointer-7}{}\fi
\noindent\textit{Source reference.} Original lemma 02JO: focused-seeds/AGP/sources/stacks-snapshot-20261011/tags.txt, source lines 3343--3343. Navigation pointer to a statement the counted proof refers to but does not contain; nothing load-bearing is delegated to it.\par\medskip
\ifdefined\hypertarget\hypertarget{sourcepointer-8}{}\fi
\noindent\textit{Source reference.} Original lemma 0538: focused-seeds/AGP/sources/stacks-snapshot-20261011/tags.txt, source lines 6596--6596. Navigation pointer to a statement the counted proof refers to but does not contain; nothing load-bearing is delegated to it.\par\medskip
\begin{lemma}
\label{lemma-flatten-on-affine-blowup}
Let $(R, \mathfrak m)$ be a local domain with fraction field $K$.
Let $S$ be a finite type $R$-algebra.
Let $M$ be a finite $S$-module.
For every valuation ring $A \subset K$ dominating $R$
there exists an ideal $I \subset \mathfrak m$ and a nonzero
element $a \in I$ such that
\begin{enumerate}
\item $I$ is finitely generated,
\item $A$ has center on $R[\frac{I}{a}]$,
\item the fibre ring of $R \to R[\frac{I}{a}]$ at $\mathfrak m$
is not zero, and
\item the strict transform $S_{I, a}$ of $S$ along $R \to R[\frac{I}{a}]$
is flat and of finite presentation over $R$, and the strict transform
$M_{I, a}$ of $M$ along $R \to R[\frac{I}{a}]$ is flat over $R$ and
finitely presented over $S_{I, a}$.
\end{enumerate}
\end{lemma}

\begin{proof}
Write $S = R[x_1, \ldots, x_n]/J$ and denote $N = S \oplus M$
viewed as a module over $P = R[x_1, \ldots, x_n]$. If we can prove the
lemma in case $S$ is a polynomial algebra over $R$, then we can
find $I, a$ satisfying (1), (2), (3) such that the strict transform
$N_{I, a}$ of $N$ along $R \to R[\frac{I}{a}]$ is flat over
$R$ and finitely presented as a module over the strict transform $P_{I, a}$
of $P$. Since $P_{I, a} = R[\frac{I}{a}][x_1, \ldots, x_n]$
(small detail omitted)
we find that the summand $S_{I, a} \subset N_{I, a}$ is flat over $R$
and finitely presented as a module over $R[\frac{I}{a}][x_1, \ldots, x_n]$.
Hence $S_{I, a}$ is finitely presented as an $R[\frac{I}{a}]$-algebra.
Moreover, the summand $M_{I, a} \subset N_{I, a}$ is flat over $R$
and finitely presented as a module over $P_{I, a}$ hence also
finitely presented as a module over $S_{I, a}$, see
Algebra, Lemma \ref{algebra-lemma-finitely-presented-over-subring}.
This reduces us to the case discussed in the next paragraph.

\medskip\noindent
Assume $S = R[x_1, \ldots, x_n]$. Choose a presentation
$$
0 \to K \to S^{\oplus r} \to M \to 0.
$$
Let $M_A$ be the quotient of $M \otimes_R A$ by its torsion submodule, see
Lemma \ref{lemma-torsion}. Then $M_A$ is a finite module over
$S_A = A[x_1, \ldots, x_n]$. By
Lemma \ref{lemma-valuation-ring-torsion-free-flat}
we see that $M_A$ is flat over $A$. By
Lemma \ref{lemma-flat-finite-type-valuation-ring-finite-presentation}
we see that $M_A$ is finitely presented. Hence there exist finitely many
elements $k_1, \ldots, k_t \in S_A^{\oplus r}$ which generate the
kernel of the presentation $S_A^{\oplus r} \to M_A$ as
an $S_A$-module. For any choice of $a \in I \subset \mathfrak m$
satisfying (1), (2), and (3) we denote $M_{I, a}$ the strict transform of
$M$ along $R \to R[\frac{I}{a}]$. It is a finite module over
$S_{I, a} = R[\frac{I}{a}][x_1, \ldots, x_n]$. By
Algebra, Lemma \ref{algebra-lemma-valuation-ring-colimit-affine-blowups}
we have $A = \colim_{I, a} R[\frac{I}{a}]$.
This implies that $S_A = \colim S_{I, a}$ and
$$
\colim M \otimes_R R[\textstyle{\frac{I}{a}}] = M \otimes_R A
$$
Choose $I, a$ and lifts $k_1, \ldots, k_t \in S_{I, a}^{\oplus r}$.
Since $M_A$ is the quotient of $M \otimes_R A$
by torsion, we see that the images of $k_1, \ldots, k_t$ in $M \otimes_R A$
are annihilated by a nonzero element $\alpha \in A$.
After replacing $I, a$ by a different pair
(recall that the colimit is filtered),
we may assume $\alpha = x/a^n$ for some $x \in I^n$ nonzero.
Then we find that $x k_1, \ldots, x k_t$ map to zero in $M \otimes_R A$.
Hence after replacing $I, a$ by a different pair
we may assume $x k_1, \ldots, x k_t$ map to zero in
$M \otimes_R R[\frac{I}{a}]$ for some nonzero $x \in R$.
Then finally replacing $I, a$ by $xI, xa$ we find that we may assume
$k_1, \ldots, k_t$ map to $a$-power torsion elements of
$M \otimes_R R[\frac{I}{a}]$. For any such pair $(I, a)$ we set
$$
M'_{I, a} = S_{I, a}^{\oplus r}/ \sum S_{I, a}k_j.
$$
Since $M_A = S_A^{\oplus r}/ \sum S_Ak_j$ we see that
$M_A = \colim_{I, a} M'_{I, a}$. At this point we finally apply
Algebra, Lemma \ref{algebra-lemma-flat-finite-presentation-limit-flat} (3)
to conclude that $M'_{I, a}$ is flat for some pair $(I, a)$ as above.
This lemma does not apply a priori to the system of strict transforms
$$
M_{I, a} = (M \otimes_R R[\textstyle{\frac{I}{a}}])/a\text{-power torsion}
$$
as the transition maps may not satisfy the assumptions of the lemma.
But now, since flatness implies torsion free
(Lemma \ref{lemma-flat-torsion-free}) and since $M_{I, a}$
is the quotient of $M'_{I, a}$ (because we arranged it so the
elements $k_1, \ldots, k_t$ map to zero in $M_{I, a}$)
by the $a$-power torsion submodule
we also conclude that $M'_{I, a} = M_{I, a}$ for such a pair and we win.
\end{proof}

\end{document}
